\newpage
\section{Ricostruzione mesh}
\label{sec:ricostruzione_mesh}

In questa fase è necessario triangolare le facce estratte, calcolare gli attributi di ciascun vertice (posizione, normale e coordinate texture) ed estrudere le facce per ottenere infine una mesh tridimensionale.

Gli algoritmi di triangolazione sono stati discussi nella Sezione \ref{sec:algoritmi_di_triangolazione}; in particolare è stata utilizzata la Constrained Delaunay Triangulation (CDT), tramite la libreria poly2tri \cite{poly2tri}. Presa in input una faccia poligonale semplice, la libreria restituisce un vettore di triangoli. Per ciascun vertice di ciascun triangolo è quindi necessario calcolare posizione, normale e coordinate texture (Figura~\ref{fig:vertex-uml}):
\begin{figure}[H]
	\centering
	\begin{tikzpicture}
		\begin{class}[text width=5.5cm]{Vertex}{0,0}
			\attribute{+ position : Vec3}
			\attribute{+ normal : Vec3}
			\attribute{+ texCoord : Vec2}
		\end{class}
	\end{tikzpicture}
	\caption{Diagramma UML della struttura \texttt{Vertex}, che rappresenta un vertice della mesh con posizione, normale e coordinate texture.}
	\label{fig:vertex-uml}
\end{figure}

L'estrusione aggiunge profondità a una regione o faccia planare selezionata, lungo una direzione prestabilita. In pratica, occorre generare i vertici di una faccia ad altezze diverse e creare dei semplici quad verticali.

Un quad ha sempre quattro vertici: in basso a sinistra (\emph{Bottom-Left}, BL), in basso a destra (\emph{Bottom-Right}, BR), in alto a sinistra (\emph{Top-Left}, TL) e in alto a destra (\emph{Top-Right}, TR); di conseguenza, viene sempre suddiviso in esattamente due triangoli. Per ciascun vertice occorre calcolare il vettore normale, perpendicolare alla superficie del quad: l'ordine in cui si percorrono i vertici per formare i triangoli (\emph{winding order}) ne determina il verso. Percorrendoli in senso antiorario (CCW) si ottiene una normale rivolta verso l'esterno, mentre percorrendoli in senso orario (CW) si ottiene una normale rivolta verso l'interno.

Le operazioni di triangolazione ed estrusione vengono applicate in modo diverso a seconda del tipo di faccia (pavimento, soffitto, muro, porta o finestra). Nel seguito, l'altezza del pavimento è indicata con $0$ e quella del soffitto con $H$.

\subsection{Pavimento e soffitto}

Il procedimento per costruire il pavimento (e analogamente il soffitto) è semplice: si calcola il bounding box 2D dell'intera casa e si definisce un'unica faccia, grande quanto l'intera abitazione, posizionata all'altezza del pavimento ($0$), quindi triangolata. La normale è rivolta verso l'alto e le coordinate texture sono definite nell'intervallo $[0,1]$.

Per il soffitto, la stessa faccia viene duplicata e traslata all'altezza $H$, con la normale rivolta verso il basso.

\subsection{Porte}

Una porta viene rappresentata come un muro pieno a partire dall'altezza dello stipite $s$ fino al soffitto, ad altezza $H$; al di sotto dello stipite, l'apertura resta libera.

In questo caso, a differenza del muro pieno, è necessario applicare la triangolazione alla faccia orizzontale posta all'altezza dello stipite $s$, per evitare che vi si formi un buco visibile. Come per i muri, non è invece necessario triangolare le facce poste ad altezza pavimento né a quella posta ad altezza soffitto.

\subsection{Finestre}

La finestra è rappresentata come un'apertura al centro del muro: si distinguono quindi una porzione di muro inferiore, dal pavimento fino al davanzale (altezza $d$), e una porzione di muro superiore, dall'architrave (altezza $a$) fino al soffitto.

Sono pertanto necessarie due estrusioni distinte: una dall'altezza $0$ all'altezza $d$, l'altra dall'altezza $a$ all'altezza $H$; a differenza dei muri pieni, è inoltre necessario triangolare le facce orizzontali poste all'altezza del davanzale e dell'architrave. Come nei casi precedenti, non è invece necessario triangolare le facce poste ad altezza pavimento e ad altezza soffitto.